\documentclass[11pt,a4paper]{article}

\usepackage[utf8]{inputenc}
\usepackage[T1]{fontenc}
\usepackage[french]{babel}
\usepackage{geometry}
\geometry{a4paper,left=2.2cm,right=2.2cm,top=2.3cm,bottom=2.3cm}
\usepackage{amsmath,amssymb}
\usepackage{booktabs}
\usepackage{array}
\usepackage{tabularx}
\usepackage{enumitem}
\usepackage{parskip}
\usepackage{fancyhdr}
\usepackage{needspace}
\usepackage[hidelinks]{hyperref}
\usepackage{url}

\let\sectionorig\section
\renewcommand{\section}{\Needspace{5\baselineskip}\sectionorig}

\newcommand{\kl}{K2/\textbackslash{}Labs}
\newcommand{\code}[1]{\texttt{\detokenize{#1}}}

\pagestyle{fancy}
\fancyhf{}
\renewcommand{\headrulewidth}{0.4pt}
\fancyhead[L]{\small Composantes du projet --- réponses}
\fancyhead[R]{\small Équipe \kl}
\fancyfoot[C]{\small\thepage}

\setlist[itemize]{topsep=2pt,itemsep=1.5pt,parsep=0pt,leftmargin=1.2em}
\setlist[enumerate]{topsep=2pt,itemsep=1.5pt,parsep=0pt,leftmargin=1.6em}
\newcolumntype{L}[1]{>{\raggedright\arraybackslash}p{#1}}

\begin{document}

\begin{titlepage}
\thispagestyle{empty}
\centering
\vspace*{2.4cm}

{\Huge\bfseries Composantes du projet}\\[0.45cm]
{\Large R\'eponses au formulaire de soumission}\\[0.3cm]
{\large GlobalAstro Health}\\[1.6cm]

\rule{0.72\textwidth}{0.6pt}\\[0.45cm]
{\large NASA Space Apps Challenge 2026}\\[0.15cm]
{\normalsize D\'efi : Create Health Monitoring Software for Astronauts on Space Missions}\\[0.45cm]
\rule{0.72\textwidth}{0.6pt}\\[1.6cm]

{\large\bfseries \'Equipe \kl}\\[0.25cm]
{\normalsize Six membres}\\[2.0cm]

\begin{tabular}{@{}ll@{}}
\small Version & \small 1.0 \\[2pt]
\small Date & \small 15 novembre 2026 \\[2pt]
\small Application & \small \url{https://www.globalastro.health} \\
\end{tabular}

\vfill
{\footnotesize Ce document contient le texte \`a reporter dans le formulaire, point par point, dans l'ordre o\`u il les demande.}
\end{titlepage}

\tableofcontents
\newpage

% ================================================================
\section{Nom du projet et r\'esum\'e g\'en\'eral}

\textbf{Nom du projet :} GlobalAstro Health

\vspace{0.25cm}
\noindent\rule{\textwidth}{0.4pt}
\vspace{0.15cm}

\noindent\emph{R\'esum\'e (158 mots, contrainte : 50 \`a 200) :}

\begin{quote}
En mission longue, personne ne peut surveiller l'\'equipage : il n'y a pas de m\'edecin \`a bord, et
le signal met des dizaines de minutes \`a atteindre la Terre. Les astronautes portent donc eux-m\^emes
la responsabilit\'e de rep\'erer les changements de leur sant\'e. Un tableau de bord qui d\'efile ne
r\'epond pas \`a cette question. GlobalAstro Health est un syst\`eme d'aide \`a la d\'ecision
construit autour d'une seule question : quand un r\'esultat justifie-t-il l'attention de l'\'equipage
ou des chirurgiens de vol ? Il \'equipe les deux extr\'emit\'es de la mission. \`A bord, sept \'ecrans
donnent \`a l'\'equipage les moyens de rep\'erer seul une d\'erive : indicateurs, os et c\oe{}ur,
sant\'e mentale, rayonnement, alertes, mode hors ligne, consignes re\c{c}ues. Au sol, quatre \'ecrans
donnent au m\'edecin l'\'etat de l'\'equipage et le cycle complet d'une consigne, de sa pr\'eparation
\`a son accus\'e de r\'eception. Les cinq dangers document\'es par le programme de recherche humaine
de la NASA structurent l'interface.
\end{quote}

\noindent\rule{\textwidth}{0.4pt}

% ================================================================
\section{Lien vers le projet et lien vers la d\'emonstration}

\textbf{Attention : le formulaire ne demande qu'UN SEUL lien} --- soit un diaporama de sept
diapositives au maximum (la page de titre compte), soit une vid\'eo de trente secondes sous-titr\'ee
en anglais. \emph{C'est le point \`a trancher avant de soumettre.}

\vspace{0.2cm}
\textbf{Lien du projet :} \url{https://www.globalastro.health}

\vspace{0.2cm}
\textbf{Lien \`a donner comme d\'emonstration --- choisir l'un des deux :}

\begin{tabularx}{\textwidth}{@{}L{3.6cm} X@{}}
\toprule
\textbf{Option} & \textbf{Ce qu'il faut faire} \\
\midrule
Diaporama (conseill\'e) & D\'eposer le diaporama de sept diapositives sur un espace public (Drive ou d\'ep\^ot de code) et coller le lien. \\
Vid\'eo de 30 s & Enregistrer l'\'ecran, sous-titres anglais obligatoires, d\'eposer sur YouTube, lien public. \\
\bottomrule
\end{tabularx}

\vspace{0.25cm}
\textbf{Trois v\'erifications avant de coller un lien :}

\begin{enumerate}
  \item ouvrir le lien dans une fen\^etre de navigation priv\'ee : aucun mot de passe, aucune connexion,
        aucune demande d'autorisation ne doit appara\^itre ;
  \item v\'erifier que le lien est accessible sans compte (ni Google, ni d\'ep\^ot) ;
  \item v\'erifier qu'aucune personne de moins de dix-huit ans n'appara\^it \`a l'image ou au son.
\end{enumerate}

\vspace{0.2cm}
\textbf{Comptes d'acc\`es \`a indiquer aux \'evaluateurs} (l'application demande une connexion et est
livr\'ee avec deux comptes pr\^ets) :

\begin{center}
\begin{tabular}{@{}L{4.4cm} L{4.4cm} L{3.0cm}@{}}
\toprule
\textbf{R\^ole} & \textbf{Identifiant} & \textbf{Mot de passe} \\
\midrule
Astronaute \`a bord & \code{cdr.mbaye} & \code{orbite2026} \\
M\'edecin au sol & \code{dr.anderson} & \code{orbite2026} \\
\bottomrule
\end{tabular}
\end{center}

\noindent{\small Un \'evaluateur peut aussi cr\'eer son propre compte depuis l'\'ecran de connexion, en
choisissant son r\^ole.}
% ================================================================
\section{D\'etails du projet}

\subsection{Que fait-il exactement ?}

Le syst\`eme sert deux m\'etiers, avec deux interfaces enti\`erement diff\'erentes selon qui se
connecte.

\vspace{0.2cm}
\textbf{\`A bord, pour l'\'equipage --- sept \'ecrans :}

\begin{enumerate}
  \item \textbf{Tableau de bord} --- les indicateurs de l'\'equipage r\'eunis, chacun avec sa
        r\'ef\'erence pr\'e-vol et sa tendance, pour rep\'erer une d\'erive et non lire un chiffre.
  \item \textbf{Os et c\oe{}ur} --- densit\'e min\'erale osseuse et param\`etres cardiovasculaires, avec
        les contre-mesures d'exercice r\'eellement utilis\'ees \`a bord de la Station spatiale.
  \item \textbf{Sant\'e mentale} --- sommeil, charge de travail et moral : le danger
        \og isolement et confinement \fg{}.
  \item \textbf{Rayonnement} --- l'exposition en \textbf{pourcentage de la limite de carri\`ere},
        jamais en millisieverts absolus, avec la contre-mesure d'abri.
  \item \textbf{Alertes} --- l'\'ecart constat\'e, la conduite \`a tenir, et l'action corrective \`a
        cocher.
  \item \textbf{Hors ligne} --- ce qui continue de fonctionner sans liaison, et la pr\'eparation de la
        demande de synchronisation.
  \item \textbf{Instructions} --- les consignes venues du sol, avec accus\'e de r\'eception.
\end{enumerate}

\vspace{0.2cm}
\textbf{Au sol, pour le m\'edecin --- quatre vues :}

\begin{enumerate}
  \item \textbf{\'Equipage} --- les membres suivis, leurs valeurs re\c{c}ues \`a la derni\`ere
        fen\^etre, l'\'ecart \`a la r\'ef\'erence ; toute personne inscrite appara\^it, y compris
        depuis un autre appareil.
  \item \textbf{Alertes} --- d\'eriv\'ees automatiquement des mesures hors plage, jamais saisies \`a
        la main.
  \item \textbf{Journal} --- les derni\`eres mesures re\c{c}ues du bord, tri\'ees par gravit\'e, avec la
        plage attendue.
  \item \textbf{Consignes} --- l'\'ecriture d'une consigne et son cycle : \emph{en attente},
        \emph{livr\'ee au bord}, \emph{lue par l'\'equipage}.
\end{enumerate}

\subsection{Comment fonctionne-t-il ?}

\textbf{\`A bord}, quinze capteurs (oxym\`etre MAX30102, thermom\`etre MAX30205, capteur de CO\textsubscript{2}
MH-Z19, centrale inertielle MPU6050, DEXA, dynamom\`etre Jamar, tensiom\`etre Omron, r\'esistance
cutan\'ee, dosim\`etre, bracelet, ceinture, laboratoire, calcul, InBody, balance) produisent vingt-cinq
param\`etres physiologiques et environnementaux. Ils alimentent un \textbf{module central \'ecrit en
Python}, sans aucune d\'ependance externe, qui calcule les seuils, attribue un statut (nominal, \`a
suivre, critique) et tient l'historique local. \textbf{Rien dans cette cha\^ine ne d\'epend du
r\'eseau} : une coupure de liaison ne doit pas interrompre la surveillance.

\textbf{La liaison est le c\oe{}ur du raisonnement.} Il n'y a pas de flux continu entre le bord et la
Terre. Les mesures restent \`a bord et partent \`a la \textbf{fen\^etre de communication}. C'est pour
cela que l'interface du sol ne propose jamais d'agir en direct : elle pr\'epare, l'\'equipage applique,
et l'accus\'e de r\'eception remonte \`a la fen\^etre suivante. Le d\'elai n'est pas un d\'etail
d'impl\'ementation, c'est ce qui d\'efinit les deux m\'etiers.

\textbf{Au sol}, une base h\'eberg\'ee (PostgreSQL) porte les comptes, les consignes et les mesures
re\c{c}ues. Le site \'etant statique, c'est son interface REST qui tient lieu de serveur applicatif.

\subsection{Quels avantages apporte-t-il ?}

\begin{itemize}
  \item \textbf{Il r\'epond \`a la question du d\'efi, pas \`a une autre.} La difficult\'e n'est pas de
        mesurer, c'est de d\'ecider quand une mesure m\'erite une action.
  \item \textbf{Il tient compte du d\'elai r\'eel.} Aucune promesse d'intervention imm\'ediate : le sol
        pr\'epare.
  \item \textbf{Il est utilisable par un non-m\'edecin.} Les \'ecrans parlent de conduite \`a tenir
        avant de parler de pathologie.
  \item \textbf{Il fonctionne hors ligne}, ce qui est la condition de survie de tout outil embarqu\'e.
  \item \textbf{Il ne ment jamais sur ses donn\'ees.} Chaque valeur affich\'ee indique sa provenance, et
        le syst\`eme dit clairement ce qui est mesur\'e, ce qui est simul\'e et ce qui reste \`a
        construire.
  \item \textbf{Il relie r\'eellement les deux extr\'emit\'es.} Une consigne \'ecrite au sol arrive \`a
        bord, et l'accus\'e de r\'eception revient : la boucle compl\`ete est d\'emontr\'ee, pas
        d\'ecrite.
\end{itemize}

% ================================================================
\section{Utilisation de l'intelligence artificielle}

\noindent{\small Le formulaire pr\'ecise que cette r\'eponse sert uniquement \`a la collecte de
donn\'ees et que les juges ne la voient pas. En revanche, ils peuvent voir tout indicateur de
g\'en\'eration par intelligence artificielle pr\'esent dans le projet. Nous avons donc choisi de
l'\'ecrire aussi dans la description du projet, plut\^ot que de le laisser deviner.}

\vspace{0.25cm}
\begin{tabularx}{\textwidth}{@{}L{3.9cm} X@{}}
\toprule
\textbf{Outil} & \textbf{O\`u et comment il a \'et\'e utilis\'e} \\
\midrule
Agent de d\'eveloppement assist\'e par un grand mod\`ele de langage & \'Ecriture et relecture du code des deux interfaces, du module central et de la base du sol ; r\'edaction de la documentation ; \'ecriture des scripts de v\'erification ; analyse des mesures pour v\'erifier la coh\'erence entre les deux postes. Utilis\'e comme outil de d\'eveloppement, sous d\'ecision et relecture humaines. \\
\textbf{\`a compl\'eter} & Documents, pr\'esentation et visuels : \`a pr\'eciser selon ce que chaque membre de l'\'equipe a utilis\'e. \\
\bottomrule
\end{tabularx}

\vspace{0.25cm}
\textbf{Pr\'ecisions que nous apportons :}

\begin{itemize}
  \item \textbf{Aucun \'el\'ement de marque ou de design de la NASA n'a \'et\'e g\'en\'er\'e ni
        utilis\'e} (logo, drapeau, identifiant de mission, police, habillage). La NASA n'approuve
        aucune entit\'e non gouvernementale am\'ericaine : nous citons ses ressources, nous ne les
        republions pas, et nous ne nous pr\'esentons pas comme approuv\'es par elle.
  \item \textbf{Les valeurs affich\'ees sont simul\'ees}, et l'interface l'indique. Aucune donn\'ee
        m\'edicale r\'eelle n'est utilis\'ee : aucune donn\'ee personnelle n'a donc transity\'e par un
        outil d'intelligence artificielle.
  \item \textbf{Aucune image de synth\`ese ne figure dans le produit.} L'identit\'e visuelle a \'et\'e
        cr\'e\'ee par l'\'equipe.
\end{itemize}
% ================================================================
\section{Donn\'ees de la NASA utilis\'ees}

\noindent{\small Rappel du r\`eglement : au moins une ressource NASA du d\'efi est obligatoire pour
concourir aux prix mondiaux.}

\vspace{0.25cm}
\begin{tabularx}{\textwidth}{@{}L{4.9cm} X@{}}
\toprule
\textbf{Ressource} & \textbf{Ce que nous en avons tir\'e, et lien} \\
\midrule
Taxonomie des cinq dangers des vols spatiaux habit\'es --- programme de recherche humaine de la NASA & La colonne vert\'ebrale de l'interface : les cinq dangers structurent les \'ecrans, et chaque \'ecran renvoie au danger correspondant. \newline \url{https://www.nasa.gov/hrp/hazards/} \\
Human Research Roadmap --- risque 77 & \emph{Risk of Bone Fracture due to Spaceflight-induced Changes to Bone} : fonde l'\'ecran \og Os et c\oe{}ur \fg{}. \newline \url{https://humanresearchroadmap.nasa.gov/risks/} \\
Human Research Roadmap --- risque 179 & \emph{Risk of Cardiovascular Adaptations Contributing to Adverse Mission Performance and Health Outcomes}. \newline \url{https://humanresearchroadmap.nasa.gov/risks/} \\
Human Research Roadmap --- risque 99 & \emph{Risk of Adverse Cognitive or Behavioral Changes and Psychiatric Disorders} : fonde l'\'ecran \og Sant\'e mentale \fg{}. \newline \url{https://humanresearchroadmap.nasa.gov/risks/} \\
Human Research Roadmap --- risque 175 & \emph{Risk of Adverse Mission Outcomes Due to Earth Independent Human-Systems Operations} : fonde l'\'ecran \og Hors ligne \fg{}. \newline \url{https://humanresearchroadmap.nasa.gov/risks/} \\
NASA-STD-3001, volume 1, \S4.2.10 & La limite de dose de carri\`ere s'exprime en risque, non en millisieverts absolus : c'est pourquoi l'\'ecran \og Rayonnement \fg{} affiche un pourcentage de cette limite. \newline \url{https://standards.nasa.gov/standard/nasa/nasa-std-3001} \\
NASA --- risque cardiovasculaire & La contre-mesure de pression n\'egative des membres inf\'erieurs. \newline \url{https://www.nasa.gov/directorates/esdmd/hhp/cv-risk/} \\
NASA --- risque musculaire & L'effet d'un programme d'exercice suivi sur la r\'ecup\'eration musculaire. \newline \url{https://www.nasa.gov/directorates/esdmd/hhp/muscle-risk/} \\
Liste des risques du programme de recherche humaine & La liste compl\`ete des risques, parcourue pour v\'erifier qu'aucun danger n'est oubli\'e. \newline \url{https://www.nasa.gov/hrp/risks/} \\
\bottomrule
\end{tabularx}

\vspace{0.25cm}
\textbf{Comment ces ressources sont utilis\'ees, concr\`etement :} elles ne sont pas cit\'ees en
bibliographie, elles \textbf{d\'ecident de la structure du produit}. Les cinq dangers d\'eterminent les
cinq familles d'\'ecrans \`a bord ; la formulation de la limite de dose par la NASA interdit
d'afficher une dose en millisieverts absolus ; la contre-mesure d'abri lors d'une \'eruption solaire
devient une action propos\'ee \`a l'\'equipage.

% ================================================================
\section{Donn\'ees des agences spatiales partenaires et autres ressources}

\noindent{\small Facultatif, mais encourag\'e.}

\vspace{0.25cm}
\begin{tabularx}{\textwidth}{@{}L{4.9cm} X@{}}
\toprule
\textbf{Ressource} & \textbf{Ce que nous en avons tir\'e, et lien} \\
\midrule
Agence spatiale canadienne --- \emph{How does radiation affect the human body in space?} & Les trois composantes du rayonnement spatial (rayons cosmiques galactiques, particules solaires, rayonnement pi\'eg\'e), les effets \`a long terme (cataractes, cancers, atteinte de la descendance), et le fait qu'en orbite basse l'atmosph\`ere et le champ magn\'etique terrestre prot\`egent encore en partie. \newline \url{https://asc-csa.gc.ca/eng/astronauts/space-medicine/radiation.asp} \\
Agence spatiale canadienne --- exp\'erience \textbf{Radi-N2} & La part des neutrons (environ 30\,\% de la dose re\c{c}ue \`a bord de la Station spatiale) et la contre-mesure : cartographier les zones les plus prot\'eg\'ees du v\'ehicule pour y replier l'\'equipage lors d'une \'eruption solaire majeure. \newline \url{https://asc-csa.gc.ca/eng/sciences/radi-n2.asp} \\
\bottomrule
\end{tabularx}

\vspace{0.25cm}
\textbf{Autres ressources techniques utilis\'ees :}

\begin{tabularx}{\textwidth}{@{}L{5.2cm} X@{}}
\toprule
\textbf{Ressource} & \textbf{Usage} \\
\midrule
React 19, TypeScript, Vite & Interfaces des deux postes \\
Python 3 (biblioth\`eque standard seule) & Module central \`a bord, sans aucune d\'ependance externe \\
PostgreSQL (service h\'eberg\'e) & Base du sol : comptes, consignes, mesures re\c{c}ues \\
Documentation des capteurs (MAX30102, MAX30205, MH-Z19, MPU6050, DEXA, Jamar, Omron) & R\'ef\'erences r\'eelles des quinze capteurs d\'ecrits \\
\bottomrule
\end{tabularx}

% ================================================================
\section{R\'ecapitulatif avant de soumettre}

\begin{tabularx}{\textwidth}{@{}L{7.4cm} X@{}}
\toprule
\textbf{Point du formulaire} & \textbf{\'Etat} \\
\midrule
Nom du projet & Fait \\
R\'esum\'e g\'en\'eral (50 \`a 200 mots) & Fait --- 158 mots \\
Lien du projet & Fait --- \url{https://www.globalastro.health} \\
Lien de la d\'emonstration (diaporama ou vid\'eo) & \textbf{\`a h\'eberger, puis \`a coller} \\
D\'etails du projet & Fait \\
Utilisation de l'intelligence artificielle & Fait --- reste \`a compl\'eter par les autres membres \\
Donn\'ees de la NASA & Fait --- neuf ressources, dont cinq risques du Human Research Roadmap \\
Donn\'ees des agences partenaires & Fait --- deux ressources de l'Agence spatiale canadienne \\
V\'erification en navigation priv\'ee & \`a faire le jour de la soumission \\
\bottomrule
\end{tabularx}

\vspace{0.6cm}
\noindent\rule{\textwidth}{0.4pt}
\vspace{0.25cm}

\noindent{\small NB : les ressources cit\'ees appartiennent \`a la NASA et \`a l'Agence spatiale
canadienne. Elles sont r\'ef\'erenc\'ees, jamais republi\'ees. La NASA n'approuve aucune entit\'e non
gouvernementale am\'ericaine et d\'ecline toute responsabilit\'e quant aux informations figurant sur
les sites non gouvernementaux.}

\end{document}
